\documentclass[hidelinks]{article}
% Page Layout
\usepackage[a4paper, total={7in, 10in}]{geometry}
% Colors
\usepackage[dvipsnames]{xcolor}
% Math
\usepackage{amsmath,amssymb,mathrsfs}
% Lists
\usepackage{enumitem}
% Section Styling
\usepackage{titlesec}
% Hyperlinks
\usepackage[unicode]{hyperref}
\usepackage[all]{hypcap}
% Code listings
\usepackage{listings}
\usepackage{inconsolata}

% --------------------------------------------------
% Color Palette
% --------------------------------------------------
\definecolor{carminered}{rgb}{1.0,0.0,0.22}
\definecolor{capri}{rgb}{0.0,0.75,1.0}
\definecolor{brightlavender}{rgb}{0.75,0.58,0.89}
\definecolor{nudelavender}{rgb}{0.78, 0.66, 0.70}
\definecolor{codebg}{rgb}{0.10,0.10,0.14}
\definecolor{codecomment}{rgb}{0.5,0.55,0.5}
\definecolor{codekeyword}{rgb}{0.55,0.78,0.98}
\definecolor{codestring}{rgb}{0.8,0.65,0.4}
\definecolor{codenumber}{rgb}{0.72,0.53,0.86}

% --------------------------------------------------
% Dark Theme
% --------------------------------------------------
\pagecolor{black}
\color{white}

% --------------------------------------------------
% Code Listing Style
% --------------------------------------------------
\lstdefinestyle{darkcode}{
  backgroundcolor=\color{codebg},
  commentstyle=\color{codecomment}\itshape,
  keywordstyle=\color{codekeyword}\bfseries,
  stringstyle=\color{codestring},
  numberstyle=\tiny\color{codenumber},
  basicstyle=\ttfamily\small\color{white},
  breakatwhitespace=false,
  breaklines=true,
  captionpos=b,
  keepspaces=true,
  numbers=left,
  numbersep=6pt,
  showspaces=false,
  showstringspaces=false,
  showtabs=false,
  tabsize=2,
  frame=single,
  rulecolor=\color{brightlavender},
  framesep=4pt,
}
\lstset{style=darkcode}
\lstdefinelanguage{CSS}{
  keywords={color,background,margin,padding,font,etc.},
  sensitive=true,
  morecomment=[l]{//},
  morecomment=[s]{/*}{*/},
  morestring=[b]",
  morestring=[b]',
}

% --------------------------------------------------
% Section Style
% --------------------------------------------------
\titleformat{\section}
  {\bfseries\Large\color{brightlavender}}
  {\thesection}{0.5em}{}

\titleformat{\subsection}
  {\bfseries\color{Thistle}}
  {\thesubsection}{0.5em}{}

\titleformat{\subsubsection}
  {\bfseries\color{nudelavender}}
  {\thesubsubsection}{0.5em}{}

% --------------------------------------------------
% Bullet Style
% --------------------------------------------------
\setlist[itemize,1]{label=$\bullet$}
\setlist[itemize,2]{label=$\circ$}
\setlist[itemize,3]{label=$\bullet$}
\setlist[itemize,4]{label=$\circ$}

% --------------------------------------------------
% Title
% --------------------------------------------------
\title{\color{carminered}\textbf{Web Programming}}
\author{\color{Thistle}Kiarash Sanei}
\date{\color{SkyBlue}June, 2026}

\begin{document}
\maketitle
\begin{Large}
\tableofcontents
\end{Large}
\newpage

% ==================================================
\section{Introduction to Web Programming}
% ==================================================

\subsection{Brief History of the Web}
\begin{itemize}
  \item The internet began as \textbf{ARPANET}, a US Department of Defense network (1960s--70s).
  \item Initial services: electronic mail and file transfer.
  \item Opened to commercial interests and most universities in the late 1980s.
  \item \textbf{WWW} (World Wide Web) created in 1989--91 by \textbf{Tim Berners-Lee}.
  \item Early browsers: Mosaic (1992), Netscape (1994), Internet Explorer (1995).
  \item Amazon.com opened in 1995; Google launched in January 1996.
\end{itemize}

\subsection{Internet vs. the WWW}
\begin{itemize}
  \item \textbf{The Internet}: A network of networks using the Internet Protocol suite to link billions of devices worldwide.
  \item Consists of millions of private, public, academic, business, and government networks.
  \item Networks are linked by electronic, wireless, and optical networking technologies.
  \item Carries services such as the WWW, email, FTP, streaming, etc.
  \item \textbf{The WWW}: A system of interlinked hypertext documents accessed via the internet using HTTP.
\end{itemize}

\subsection{Web Architecture Evolution}
\begin{itemize}
  \item \textbf{1990s}: Static websites --- simple HTML pages served directly from a web server.
  \item \textbf{Early 2000s}: Dynamic web --- server-side rendering with PHP, ASP, and JSP.
  \item \textbf{Mid-2000s}: MVC frameworks --- Rails, Django, and Spring MVC for cleaner app design.
  \item \textbf{2010s}: Single-Page Apps (SPA) --- React, Angular, and Vue move logic to the client side.
  \item \textbf{Late 2010s}: Serverless \& Microservices --- apps split into small, independent services.
  \item \textbf{2020s}: Modern hybrids --- Next.js, Remix, and Astro blend SSR, SPA, and static generation.
\end{itemize}

\subsection{The Layers of the Web}
\begin{center}
  \texttt{Client (Front-end)} $\xrightarrow{\text{HTTP Request}}$ \texttt{Network} $\xrightarrow{\text{HTTP Response}}$ \texttt{Server (Back-end)}
\end{center}
\begin{itemize}
  \item The \textbf{front-end} handles the user interface and runs in the browser.
  \item The \textbf{back-end} handles business logic, databases, and API endpoints.
  \item Communication happens over \textbf{HTTP} using URLs and API endpoints.
\end{itemize}

\subsection{The Holy Trinity of the Web}
\begin{itemize}
  \item \textbf{HTML} --- Structure (what the content \emph{is}).
  \item \textbf{CSS} --- Style (how it \emph{looks}).
  \item \textbf{JavaScript} --- Behaviour (how it \emph{reacts}).
\end{itemize}

\subsection{Technology Stacks}
\begin{itemize}
  \item A \textbf{stack} is the set of technologies used to build a web application.
  \item Combines frontend, backend, database, and infrastructure layers.
  \item Each layer handles a different part of the application (UI, logic, data, hosting).
  \item Choosing the right stack affects performance, scalability, and developer workflow.
  \item \textbf{Full-stack developers} work across all layers.
\end{itemize}

\subsection{Developer Roles}
\begin{itemize}
  \item \textbf{Front-End Developer}: Builds interactive, user-friendly web interfaces using HTML, CSS, JavaScript, and frameworks like React or Vue.
  \item \textbf{Back-End Developer}: Develops server-side logic, APIs, and database integrations using Python, Node.js, Golang, or .NET.
  \item \textbf{Full-Stack Developer}: Combines front-end and back-end skills to deliver complete web solutions.
  \item \textbf{Web DevOps \& Cloud Engineer}: Manages deployments, CI/CD pipelines, and scalable cloud hosting.
\end{itemize}

\subsection{Software Engineering Principles}
\begin{itemize}
  \item \textbf{Clean Code}: Ensures readability, consistency, and collaboration across teams.
  \item \textbf{Maintainable Code}: Easier debugging, feature additions, and long-term scalability.
  \item \textbf{Dead Code Removal}: Eliminates unused, redundant logic to reduce complexity.
  \item \textbf{Fail Fast Approach}: Detects errors early, improving reliability and lowering costs.
\end{itemize}

\subsection{LLM Best Practices}
\begin{itemize}
  \item \textbf{Define clear goals}: Know exactly what you want the model to achieve before prompting.
  \item \textbf{Prompt precisely}: Use clear, specific, and contextual instructions.
  \item \textbf{Iterate \& refine}: Test, review, and adjust prompts or parameters for consistency.
  \item \textbf{Keep context short \& relevant}: Provide only necessary information to reduce confusion.
  \item \textbf{Validate outputs}: Always fact-check and post-process model responses.
  \item \textbf{Respect safety \& ethics}: Ensure responsible AI use.
  \item Use of LLMs in this course is permitted \textbf{as long as you understand the code}.
\end{itemize}

\newpage
% ==================================================
\section{HTML --- The Building Blocks of the Web}
% ==================================================

\subsection{What is HTML?}
\begin{itemize}
  \item \textbf{HTML} stands for \textbf{HyperText Markup Language}.
  \item Used to \textbf{structure and describe} the content of a webpage.
  \item Consists of \textbf{elements} that describe different types of content.
  \item Web browsers understand HTML and render it as websites.
  \item \textbf{HTML is not a programming language} --- it is a markup language.
\end{itemize}

\subsection{HTML Document Structure}
\begin{lstlisting}[language=HTML, caption={Basic HTML5 document structure}]
<!DOCTYPE html>
<html lang="en">
  <head>
    <meta charset="UTF-8" />
    <meta name="viewport" content="width=device-width, initial-scale=1.0" />
    <title>My Page</title>
    <link rel="stylesheet" href="style.css" />
  </head>
  <body>
    <h1>Hello, World!</h1>
    <p>This is a paragraph.</p>
  </body>
</html>
\end{lstlisting}
\begin{itemize}
  \item \texttt{<!DOCTYPE html>} declares the HTML version (HTML5).
  \item \texttt{<html>} wraps the entire document.
  \item \texttt{<head>} contains meta information (not displayed directly).
  \item \texttt{<body>} holds all visible content.
  \item Tags are structured hierarchically and must be properly opened and closed.
\end{itemize}

\subsection{HTML Elements}
An HTML element typically has:
\begin{itemize}
  \item An \textbf{opening tag}: \texttt{<button>}
  \item A \textbf{closing tag}: \texttt{</button>}
  \item \textbf{Attributes} (key-value pairs): \texttt{id="submit"}
  \item \textbf{Child elements or content} inside the tags.
\end{itemize}

\subsection{The \texttt{<head>} Element}
\begin{itemize}
  \item Contains information \emph{about} the page, not displayed directly.
  \item Helps browsers, search engines, and devices understand your page.
  \item \texttt{<title>} --- sets the page title (shown on the browser tab).
  \item \texttt{<meta>} --- defines metadata such as charset, author, viewport, and description.
  \item \texttt{<link>} --- connects external resources like stylesheets or favicons.
  \item \texttt{<script>} --- links or embeds JavaScript files.
\end{itemize}

\subsection{Text Elements}
\begin{itemize}
  \item \texttt{<p>} --- paragraph text.
  \item \texttt{<h1>} through \texttt{<h6>} --- headings (largest to smallest).
  \item \texttt{<strong>} or \texttt{<b>} --- bold font (\texttt{<strong>} carries semantic emphasis).
  \item \texttt{<em>} or \texttt{<i>} --- italic font (\texttt{<em>} carries semantic emphasis).
  \item \texttt{<code>} --- displays a piece of computer code inline.
  \item \texttt{<pre>} --- preserves whitespace and displays preformatted text (often wraps \texttt{<code>}).
  \item \texttt{<span>} --- generic inline container for styling parts of text.
  \item \texttt{<blockquote>} --- block-level quotation.
  \item \texttt{<br>} --- line break (void element, no closing tag).
  \item \texttt{<hr>} --- horizontal rule / thematic break (void element, no closing tag).
\end{itemize}

\begin{lstlisting}[language=HTML, caption={Text element examples}]
<h1>Main Title</h1>
<h2>Subtitle</h2>
<p>This is a <strong>bold</strong> and <em>italic</em> paragraph.</p>
<p>Inline code: <code>console.log("hello")</code></p>
<pre><code>
function greet(name) {
  return "Hello, " + name;
}
</code></pre>
\end{lstlisting}

\subsection{Hyperlinks}
\begin{itemize}
  \item Created using the \texttt{<a>} (anchor) tag.
  \item The \texttt{href} attribute specifies the URL or location to link to.
  \item \texttt{target="\_blank"} opens a link in a new tab.
  \item The \texttt{download} attribute triggers a file download.
  \item \textbf{Internal links}: \texttt{href="/about"}
  \item \textbf{External links}: \texttt{href="https://wikipedia.org"}
  \item \textbf{Anchor links}: \texttt{href="\#section-id"} (jumps to a section on the same page)
  \item \textbf{Mail links}: \texttt{href="mailto:info@example.com"}
  \item \textbf{Tel links}: \texttt{href="tel:+9821000000"}
\end{itemize}

\begin{lstlisting}[language=HTML, caption={Hyperlink examples}]
<a href="https://google.com" target="_blank">Visit Google</a>
<a href="/contact">Contact Us</a>
<a href="#top">Back to Top</a>
<a href="mailto:hello@example.com">Email Us</a>
<a href="report.pdf" download>Download Report</a>
\end{lstlisting}

\subsection{List Elements}
\begin{itemize}
  \item \texttt{<ul>} --- unordered list (bulleted).
  \item \texttt{<ol>} --- ordered list (numbered).
  \item \texttt{<li>} --- list item (used inside both \texttt{<ul>} and \texttt{<ol>}).
  \item \texttt{<dl>}, \texttt{<dt>}, \texttt{<dd>} --- description list with terms and definitions.
\end{itemize}

\begin{lstlisting}[language=HTML, caption={List examples}]
<ul>
  <li>Apple</li>
  <li>Banana</li>
</ul>

<ol>
  <li>Step one</li>
  <li>Step two</li>
</ol>

<dl>
  <dt>HTML</dt>
  <dd>HyperText Markup Language</dd>
</dl>
\end{lstlisting}

\subsection{Media Elements}
\begin{itemize}
  \item \texttt{<img src="..." alt="...">} --- embed an image (\texttt{alt} is required for accessibility).
  \item \texttt{<audio controls>} --- embed an audio player.
  \item \texttt{<video controls width="...">} --- embed a video player.
  \item \texttt{<source>} --- specifies multiple media sources inside \texttt{<audio>} or \texttt{<video>}.
  \item \texttt{<iframe>} --- embeds another webpage or external content.
\end{itemize}

\begin{lstlisting}[language=HTML, caption={Media elements}]
<img src="photo.jpg" alt="A landscape photo" width="600" />

<audio controls>
  <source src="song.mp3" type="audio/mpeg" />
  Your browser does not support audio.
</audio>

<video controls width="640">
  <source src="clip.mp4" type="video/mp4" />
  Your browser does not support video.
</video>
\end{lstlisting}

\subsection{Tables}
\begin{itemize}
  \item \texttt{<table>} --- defines the table.
  \item \texttt{<tr>} --- table row.
  \item \texttt{<th>} --- header cell (bold by default).
  \item \texttt{<td>} --- data cell.
  \item \texttt{<thead>}, \texttt{<tbody>}, \texttt{<tfoot>} --- optional sections for organizing large tables.
  \item Use \texttt{colspan} and \texttt{rowspan} attributes to merge cells.
\end{itemize}

\begin{lstlisting}[language=HTML, caption={Table example with merged cell}]
<table border="1">
  <thead>
    <tr>
      <th>Name</th>
      <th>Age</th>
      <th colspan="2">Contact</th>
    </tr>
  </thead>
  <tbody>
    <tr>
      <td>Alice</td>
      <td>22</td>
      <td>alice@mail.com</td>
      <td>+1-555-0100</td>
    </tr>
  </tbody>
</table>
\end{lstlisting}

\subsection{Forms}
\begin{itemize}
  \item \texttt{<form>} --- wraps all form controls; has \texttt{action} (where to send) and \texttt{method} (\texttt{GET}/\texttt{POST}).
  \item \texttt{<input>} --- a variety of controls selected via the \texttt{type} attribute.
  \item \texttt{<textarea>} --- multi-line text input.
  \item \texttt{<select>} + \texttt{<option>} --- dropdown menu.
  \item \texttt{<button>} --- clickable button.
  \item \texttt{<label>} --- associates descriptive text with a control (use \texttt{for} attribute).
  \item \texttt{<fieldset>} + \texttt{<legend>} --- groups related controls.
\end{itemize}

\begin{lstlisting}[language=HTML, caption={Sample HTML form}]
<form action="/submit" method="POST">
  <label for="name">Name:</label>
  <input type="text" id="name" name="name" required />

  <label for="email">Email:</label>
  <input type="email" id="email" name="email" />

  <label for="role">Role:</label>
  <select id="role" name="role">
    <option value="dev">Developer</option>
    <option value="design">Designer</option>
  </select>

  <label>
    <input type="checkbox" name="agree" /> I agree to the terms
  </label>

  <button type="submit">Send</button>
</form>
\end{lstlisting}

\subsection{Identifying Elements: \texttt{id} and \texttt{class}}
\begin{itemize}
  \item \texttt{id} --- a \textbf{unique} identifier for a single element on the page.
  \item \texttt{class} --- groups multiple elements with shared styling or behaviour.
  \item An element can have only \textbf{one} \texttt{id}, but \textbf{multiple} class values (space-separated).
  \item \texttt{id} is used for anchors, scripting, or targeting specific elements.
  \item \texttt{class} is used mostly for styling and reusable components.
\end{itemize}

\subsection{Sectioning the Page}

\subsubsection{Using \texttt{<div>}}
\begin{itemize}
  \item A generic container to group HTML elements.
  \item Does not carry any semantic meaning on its own.
  \item Commonly used with \texttt{class} or \texttt{id} to apply styles or behaviour.
\end{itemize}

\subsubsection{Using Semantic Tags (Preferred)}
\begin{itemize}
  \item \texttt{<header>} --- introductory content or navigation links.
  \item \texttt{<nav>} --- major navigation block.
  \item \texttt{<main>} --- dominant content of the \texttt{<body>}.
  \item \texttt{<section>} --- thematic grouping of content.
  \item \texttt{<article>} --- self-contained composition (blog post, news story).
  \item \texttt{<aside>} --- tangentially related content (sidebar).
  \item \texttt{<footer>} --- footer for the page or a section.
\end{itemize}
Semantic tags improve readability, accessibility, and SEO compared to using \texttt{<div>} everywhere.

\subsection{Inline vs.\ Block Elements}
\begin{itemize}
  \item \textbf{Block elements}: Start on a new line and take up the full width available (e.g.\ \texttt{<div>}, \texttt{<p>}, \texttt{<h1>}, \texttt{<ul>}).
  \item \textbf{Inline elements}: Do not start on a new line; take only as much width as needed (e.g.\ \texttt{<span>}, \texttt{<a>}, \texttt{<strong>}, \texttt{<img>}).
  \item Inline elements \textbf{cannot} contain block-level elements.
\end{itemize}

\subsection{Comments}
\begin{lstlisting}[language=HTML, caption={HTML comment syntax}]
<!-- This is a comment and will not appear in the browser -->
<p>This text is visible.</p>
\end{lstlisting}

\subsection{Debugging HTML}
\begin{itemize}
  \item Use the browser's \textbf{Inspect Element} (DevTools) to analyze structure, debug layout, and experiment live.
  \item Always look for the \textbf{root cause} of a problem, not just the symptom.
  \item \textit{``Fix the problem, not the blame.''} --- The Pragmatic Programmer.
  \item The first rule of debugging: \textbf{Don't Panic}.
  \item Explaining a bug to someone else (``rubber duck debugging'') often reveals the issue.
\end{itemize}

\newpage
% ==================================================
\section{CSS --- Pages That Look Elegant}
% ==================================================

\subsection{What is CSS?}
\begin{itemize}
  \item \textbf{CSS} stands for \textbf{Cascading Style Sheets}.
  \item Used to \textbf{style} HTML elements --- controls colors, fonts, layout, and animations.
  \item CSS \textbf{separates content} (HTML) from \textbf{presentation} (design).
  \item Helps create responsive, modern, and user-friendly designs.
\end{itemize}

\subsection{CSS Rule Syntax}
\begin{lstlisting}[language=CSS, caption={CSS rule anatomy}]
selector {
  property: value;
}

/* Example */
p {
  color: white;
  font-size: 16px;
}
\end{lstlisting}

\subsection{CSS Selectors}
\begin{itemize}
  \item \texttt{element} --- targets all elements of a type (e.g.\ \texttt{p}, \texttt{div}).
  \item \texttt{.class} --- targets elements with a specific class.
  \item \texttt{\#id} --- targets a specific element with an ID.
  \item \texttt{*} --- targets all elements (universal selector).
  \item \texttt{A, B} --- targets both A and B elements.
  \item \texttt{A B} --- targets B that is a \textbf{descendant} of A.
  \item \texttt{A > B} --- targets B that is a \textbf{direct child} of A.
  \item \texttt{A + B} --- targets B that is the \textbf{adjacent sibling} immediately after A.
  \item \texttt{A \textasciitilde{} B} --- targets all siblings B after A.
\end{itemize}

\subsection{Pseudo-class Selectors}
\begin{itemize}
  \item \texttt{:hover} --- when the mouse is over the element.
  \item \texttt{:focus} --- when the element is focused (e.g.\ an input field).
  \item \texttt{:first-child} / \texttt{:last-child} --- first or last child of its parent.
  \item \texttt{:nth-child(n)} --- every $n$th child (e.g.\ \texttt{:nth-child(2n)} for even rows).
  \item \texttt{:not(.active)} --- all elements except those with \texttt{.active}.
  \item \texttt{:checked}, \texttt{:disabled}, \texttt{:placeholder-shown} --- form-state selectors.
\end{itemize}

\subsection{Attribute Selectors}
\begin{itemize}
  \item \texttt{[type]} --- elements that have the \texttt{type} attribute.
  \item \texttt{[type="text"]} --- elements where \texttt{type} equals \texttt{"text"}.
  \item \texttt{[href\textasciicircum="https"]} --- \texttt{href} starting with \texttt{"https"}.
  \item \texttt{[href\$=".pdf"]} --- \texttt{href} ending with \texttt{".pdf"}.
  \item \texttt{[data-role*="btn"]} --- attribute containing \texttt{"btn"}.
\end{itemize}

\subsection{Element Selectors (Type Selectors)}
\begin{itemize}
  \item \texttt{p} --- selects all paragraphs.
  \item \texttt{div} --- selects all \texttt{<div>} containers.
  \item \texttt{h1}, \texttt{h2}, \texttt{h3} --- selects all headings of a specific level.
  \item \texttt{ul li} --- descendant combination: all \texttt{<li>} inside any \texttt{<ul>}.
\end{itemize}

\subsection{Class Selectors}
\begin{itemize}
  \item \texttt{.classname} --- selects all elements that have \texttt{class="classname"}.
  \item \texttt{div.classname} --- selects only \texttt{<div>} elements with that class.
  \item \texttt{.class1.class2} --- selects elements that have \textbf{both} classes.
  \item \texttt{.parent .child} --- descendant combination with classes.
\end{itemize}

\subsection{Pseudo-element Selectors}
\begin{itemize}
  \item \texttt{::before} --- inserts content before the element’s real content (requires \texttt{content}).
  \item \texttt{::after} --- inserts content after the element’s real content (requires \texttt{content}).
  \item \texttt{::first-line} --- styles the first line of a block of text.
  \item \texttt{::first-letter} --- styles the first character of a block of text.
  \item \texttt{::selection} --- styles the part of the text that the user highlights.
  \item \texttt{::marker} --- styles the bullet or number in list items.
  \item \texttt{::placeholder} --- styles the placeholder text inside an \texttt{<input>} or \texttt{<textarea>}.
\end{itemize}

\subsection{Including CSS in HTML}
\begin{itemize}
  \item \textbf{Inline CSS}: \texttt{<p style="color:red;">Hello</p>} --- highest specificity, not recommended for large projects.
  \item \textbf{Internal CSS}: \texttt{<style>} block inside \texttt{<head>} --- good for single-page documents.
  \item \textbf{External CSS}: \texttt{<link rel="stylesheet" href="style.css">} --- best practice; styles apply to multiple pages.
\end{itemize}

\begin{lstlisting}[language=HTML, caption={Linking an external stylesheet}]
<head>
  <link rel="stylesheet" href="styles.css" />
</head>
\end{lstlisting}

\subsection{CSS Specificity and the Cascade}
Specificity determines which CSS rule wins when multiple rules target the same element.
\begin{enumerate}
  \item \textbf{Inline styles} --- highest priority.
  \item \textbf{ID selectors} (\texttt{\#id}) --- very high.
  \item \textbf{Class, pseudo-class, attribute selectors} --- medium.
  \item \textbf{Element and pseudo-element selectors} --- lowest.
\end{enumerate}
\begin{itemize}
  \item When specificity is \textbf{equal}, the \textbf{last rule declared} wins.
  \item \texttt{!important} overrides all other rules (use sparingly).
\end{itemize}

\subsection{CSS Measurement Units}
\begin{itemize}
  \item \textbf{Absolute units}: \texttt{px} (pixels), \texttt{pt} (points), \texttt{cm}, \texttt{mm}, \texttt{in}.
  \item \textbf{Relative units}:
    \begin{itemize}
      \item \texttt{em} --- relative to the \textbf{element's} font-size.
      \item \texttt{rem} --- relative to the \textbf{root} element's font-size.
      \item \texttt{vw} / \texttt{vh} --- 1\% of the viewport width / height.
      \item \texttt{\%} --- relative to the \textbf{parent} element.
      \item \texttt{ch} --- width of the ``0'' character; useful for text containers.
    \end{itemize}
\end{itemize}

\subsection{Text and Font Properties}
\begin{lstlisting}[language={}, caption={Common text properties}]
p {
  font-family: 'Inter', Arial, sans-serif;
  font-size: 1rem;         /* 16px by default */
  font-weight: 600;        /* normal | bold | 100-900 */
  line-height: 1.6;
  text-align: justify;     /* left | center | right | justify */
  text-decoration: none;   /* underline | line-through */
  letter-spacing: 0.05em;
  text-transform: uppercase;
}
\end{lstlisting}

\subsection{Color Properties}
\begin{itemize}
  \item \texttt{color} --- sets the text color.
  \item \texttt{background-color} --- sets the element background.
  \item \texttt{background-image} --- sets a background image (\texttt{url(...)}).
  \item Colors can be expressed as:
    \begin{itemize}
      \item Named: \texttt{red}, \texttt{cornflowerblue}
      \item HEX: \texttt{\#ff4455}
      \item RGB / RGBA: \texttt{rgb(255, 0, 0)} / \texttt{rgba(255, 0, 0, 0.5)}
      \item HSL / HSLA: \texttt{hsl(0, 100\%, 50\%)}
    \end{itemize}
  \item Use high contrast for readability and accessibility.
\end{itemize}

\subsection{The CSS Box Model}
Every element is a rectangular box composed of (inside-out):
\[
  \text{content} \;\to\; \text{padding} \;\to\; \text{border} \;\to\; \text{margin}
\]
\begin{itemize}
  \item \texttt{width} / \texttt{height} --- size of the content area.
  \item \texttt{padding} --- space inside the border, around content.
  \item \texttt{border} --- line surrounding the padding.
  \item \texttt{margin} --- space \textbf{outside} the border (collapses between vertical siblings).
  \item \texttt{box-sizing: border-box} --- makes \texttt{width}/\texttt{height} include padding and border (recommended).
\end{itemize}

\begin{lstlisting}[language={}, caption={Box model reset (best practice)}]
*, *::before, *::after {
  box-sizing: border-box;
  margin: 0;
  padding: 0;
}
\end{lstlisting}

\subsection{Borders and Outlines}
\begin{itemize}
  \item \texttt{border}: shorthand --- \texttt{2px solid \#fff}.
  \item Border styles: \texttt{solid}, \texttt{dashed}, \texttt{dotted}, \texttt{double}, \texttt{none}.
  \item \texttt{border-radius} --- rounds corners; \texttt{50\%} creates a circle.
  \item \texttt{outline} --- drawn outside the border; \textbf{does not affect layout} (no space taken).
\end{itemize}

\subsection{CSS Display and Layout}
\begin{itemize}
  \item \texttt{display: block} --- element starts on a new line and fills available width.
  \item \texttt{display: inline} --- flows with surrounding text.
  \item \texttt{display: inline-block} --- flows inline but respects width/height.
  \item \texttt{display: none} --- removes element from layout entirely.
  \item \texttt{display: flex} --- activates Flexbox on the container.
  \item \texttt{display: grid} --- activates CSS Grid on the container.
  \item \texttt{visibility: hidden} --- hides element but \textbf{keeps its space}.
\end{itemize}

\subsection{CSS Floats}
\begin{itemize}
  \item \texttt{float: left | right} --- element floats to the side; surrounding text wraps around it.
  \item \texttt{clear: left | right | both} --- prevents the following element from wrapping around floats.
  \item Mostly superseded by Flexbox and Grid for layout, but still used for text-wrapping around images.
\end{itemize}

\subsection{Flexbox}
\begin{itemize}
  \item A \textbf{one-dimensional} layout method (one row \emph{or} one column at a time).
  \item \texttt{display: flex} --- makes the element a flex \textbf{container}; direct children become flex \textbf{items}.
  \item \texttt{flex-direction: row | column | row-reverse | column-reverse}.
  \item \texttt{justify-content} --- aligns items along the \textbf{main axis} (e.g.\ \texttt{center}, \texttt{space-between}).
  \item \texttt{align-items} --- aligns items along the \textbf{cross axis} (e.g.\ \texttt{center}, \texttt{stretch}).
  \item \texttt{flex-wrap: wrap} --- allows items to wrap to the next line.
  \item \texttt{gap} --- spacing between flex items.
  \item On a flex item: \texttt{flex: 1} distributes available space equally; \texttt{align-self} overrides \texttt{align-items}.
\end{itemize}

\begin{lstlisting}[language={}, caption={Flexbox centering example}]
.container {
  display: flex;
  justify-content: center;
  align-items: center;
  gap: 1rem;
  flex-wrap: wrap;
}
\end{lstlisting}

\subsection{CSS Grid}
\begin{itemize}
  \item A \textbf{two-dimensional} layout method (rows \emph{and} columns simultaneously).
  \item \texttt{display: grid} --- makes the element a grid container.
  \item \texttt{grid-template-columns} --- defines column sizes (e.g.\ \texttt{repeat(3, 1fr)}, \texttt{200px auto}).
  \item \texttt{grid-template-rows} --- defines row sizes.
  \item \texttt{gap} --- gutters between rows and columns.
  \item \texttt{grid-column} / \texttt{grid-row} --- specifies where an item starts and ends.
  \item \texttt{grid-template-areas} --- assigns named areas; \texttt{grid-area} places items into them.
\end{itemize}

\begin{lstlisting}[language={}, caption={CSS Grid layout example}]
.layout {
  display: grid;
  grid-template-columns: 200px 1fr;
  grid-template-rows: auto 1fr auto;
  grid-template-areas:
    "header header"
    "sidebar main"
    "footer footer";
  gap: 1rem;
}

.header { grid-area: header; }
.sidebar { grid-area: sidebar; }
.main   { grid-area: main; }
.footer { grid-area: footer; }
\end{lstlisting}

\subsection{CSS Positioning}
\begin{itemize}
  \item \texttt{static} --- default; element flows normally in the page.
  \item \texttt{relative} --- positioned relative to its \textbf{normal location}; still in flow.
  \item \texttt{absolute} --- removed from flow; positioned relative to the nearest \textbf{positioned ancestor}.
  \item \texttt{fixed} --- removed from flow; positioned relative to the \textbf{viewport}; stays when scrolling.
  \item \texttt{sticky} --- toggles between \texttt{relative} and \texttt{fixed} based on scroll position.
\end{itemize}

\begin{lstlisting}[language={}, caption={Sticky header example}]
header {
  position: sticky;
  top: 0;
  z-index: 100;
  background-color: #111;
}
\end{lstlisting}

\subsection{Responsive Design and Media Queries}
\begin{itemize}
  \item Responsive design makes a webpage adjust to any screen size.
  \item \textbf{Media queries} apply styles based on device/viewport properties.
  \item Syntax: \texttt{@media (condition) \{ ... \}}
  \item Common conditions: \texttt{max-width}, \texttt{min-width}, \texttt{orientation}, \texttt{prefers-color-scheme}.
  \item \textbf{Mobile-first approach}: write base styles for mobile, then use \texttt{min-width} queries to enhance for larger screens.
\end{itemize}

\begin{lstlisting}[language={}, caption={Responsive breakpoints (mobile-first)}]
/* Mobile base styles */
.card { font-size: 14px; padding: 1rem; }

/* Tablet */
@media (min-width: 768px) {
  .card { font-size: 16px; padding: 1.5rem; }
}

/* Desktop */
@media (min-width: 1024px) {
  .card { font-size: 18px; padding: 2rem; }
}
\end{lstlisting}

\subsection{CSS Variables (\texttt{:root})}
\begin{itemize}
  \item Define variables on the \texttt{:root} selector to make them globally available.
  \item Variable names start with \texttt{--}.
  \item Access with \texttt{var(--variable-name)}.
  \item Makes styles consistent and easy to update globally.
\end{itemize}

\begin{lstlisting}[language={}, caption={CSS variables example}]
:root {
  --primary: #7a5af8;
  --bg: #0d0d0d;
  --text: #e8e8e8;
  --radius: 8px;
}

button {
  background-color: var(--primary);
  color: var(--text);
  border-radius: var(--radius);
}
\end{lstlisting}

\newpage
% ==================================================
\section{JavaScript Fundamentals}
% ==================================================

JavaScript is the programming language of the web --- it adds \textbf{behaviour} to HTML and CSS, and (via Node.js) runs on the server too.

\subsection{Variables and Data Types}
\begin{itemize}
  \item \texttt{var} --- function-scoped, hoisted, legacy (avoid in modern code).
  \item \texttt{let} --- block-scoped, mutable.
  \item \texttt{const} --- block-scoped, \textbf{immutable binding} (the value itself may still be mutable for objects/arrays).
  \item \textbf{Primitive types}: \texttt{number}, \texttt{string}, \texttt{boolean}, \texttt{null}, \texttt{undefined}, \texttt{symbol}, \texttt{bigint}.
  \item \textbf{Reference types}: \texttt{object}, \texttt{array}, \texttt{function}.
\end{itemize}

\begin{lstlisting}[language=Java, caption={Variable declarations}]
let age = 25;            // mutable
const name = "Alice";   // immutable binding
var legacy = true;      // avoid -- function-scoped

// typeof operator
typeof 42;          // "number"
typeof "hello";     // "string"
typeof null;        // "object" (historical quirk!)
typeof undefined;   // "undefined"
typeof [];          // "object"
typeof {};          // "object"
typeof function(){}; // "function"
\end{lstlisting}

\subsection{Type Coercion}
JavaScript is \textbf{dynamically typed} and performs implicit type conversion, which can produce surprising results.

\begin{lstlisting}[language=Java, caption={Type coercion examples}]
// Implicit coercion (loose equality ==)
"5" == 5;       // true  (string coerced to number)
0 == false;     // true
null == undefined; // true
null == 0;      // false

// Always prefer strict equality ===
"5" === 5;      // false (no coercion)
0 === false;    // false

// Addition vs concatenation
"5" + 2;        // "52"  (number coerced to string)
"5" - 2;        // 3     (string coerced to number)
"5" * "3";      // 15
+"5";           // 5     (unary + converts string to number)
!!"hello";      // true  (double negation to boolean)

// Falsy values: false, 0, "", null, undefined, NaN
// Everything else is truthy
\end{lstlisting}

\subsection{Functions}

\subsubsection{Hoisting}
\textbf{Hoisting} is JavaScript's behaviour of moving declarations to the top of their scope \emph{before} code executes.
\begin{itemize}
  \item \textbf{Function declarations} are fully hoisted --- the entire function body is available before its line in the code.
  \item \textbf{\texttt{var} variables} are hoisted but initialised to \texttt{undefined} until the assignment line is reached.
  \item \textbf{\texttt{let} and \texttt{const}} are hoisted but \textbf{not initialised} --- accessing them before their declaration throws a \texttt{ReferenceError} (the \textbf{Temporal Dead Zone}).
  \item \textbf{Function expressions} and \textbf{arrow functions} assigned to variables follow the hoisting rules of \texttt{var}/\texttt{let}/\texttt{const} and are \textbf{not} callable before their declaration.
\end{itemize}

\begin{lstlisting}[language=Java, caption={Hoisting behaviour}]
// Function declaration -- WORKS (fully hoisted)
console.log(add(2, 3)); // 5
function add(a, b) { return a + b; }

// var -- hoisted as undefined, no error but unexpected value
console.log(x); // undefined (not an error)
var x = 10;
console.log(x); // 10

// let/const -- Temporal Dead Zone (TDZ)
console.log(y); // ReferenceError: Cannot access 'y' before initialization
let y = 20;

// Function expression -- NOT hoisted
console.log(multiply(2, 3)); // TypeError: multiply is not a function
const multiply = (a, b) => a * b;
\end{lstlisting}

\subsubsection{Function Types}
\begin{lstlisting}[language=Java, caption={All function declaration styles}]
// 1. Function declaration (hoisted, has own `this`)
function add(a, b) {
  return a + b;
}

// 2. Function expression (not hoisted, has own `this`)
const subtract = function(a, b) {
  return a - b;
};

// 3. Arrow function (not hoisted, NO own `this`, NO `arguments`)
const multiply = (a, b) => a * b;

// Single parameter -- parens optional
const double = n => n * 2;

// No parameters -- parens required
const greetWorld = () => "Hello, World!";

// Multi-line arrow function -- need explicit return
const divide = (a, b) => {
  if (b === 0) throw new Error("Division by zero");
  return a / b;
};
\end{lstlisting}

\subsubsection{Parameters}
\begin{lstlisting}[language=Java, caption={Default, rest, and spread in functions}]
// Default parameters
function greet(name = "World", greeting = "Hello") {
  return `${greeting}, ${name}!`;
}
greet();              // "Hello, World!"
greet("Alice");       // "Hello, Alice!"
greet("Bob", "Hi");   // "Hi, Bob!"

// Rest parameters (must be last)
function sum(first, ...rest) {
  return rest.reduce((acc, n) => acc + n, first);
}
sum(1, 2, 3, 4); // 10

// Spread to pass array as arguments
const nums = [1, 2, 3];
Math.max(...nums); // 3
\end{lstlisting}

\subsubsection{Higher-Order Functions}
A \textbf{higher-order function} is a function that either \textbf{takes a function as an argument} or \textbf{returns a function} (or both). This is the foundation of functional patterns in JS.

\begin{lstlisting}[language=Java, caption={Higher-order functions}]
// Takes a function as argument
function repeat(n, action) {
  for (let i = 0; i < n; i++) action(i);
}
repeat(3, console.log); // prints 0, 1, 2

// Returns a function (function factory)
function makeMultiplier(factor) {
  return (n) => n * factor; // closure over `factor`
}
const triple = makeMultiplier(3);
triple(5); // 15

// Built-in HOFs: map, filter, reduce, forEach, sort
[1, 2, 3].map(n => n ** 2);       // [1, 4, 9]
[1, 2, 3, 4].filter(n => n % 2 === 0); // [2, 4]
\end{lstlisting}

\subsubsection{Immediately Invoked Function Expressions (IIFE)}
An \textbf{IIFE} is a function that is defined and called immediately. It was commonly used to create a private scope before ES modules existed.

\begin{lstlisting}[language=Java, caption={IIFE pattern}]
// Classic IIFE
(function() {
  const secret = "hidden";
  console.log("IIFE ran, secret:", secret);
})();
// `secret` is not accessible outside

// Arrow IIFE
(() => {
  console.log("Arrow IIFE");
})();

// IIFE with return value
const result = (() => {
  const x = 10, y = 20;
  return x + y;
})();
console.log(result); // 30
\end{lstlisting}

\subsection{Objects and JSON}
\begin{itemize}
  \item Objects group related data and behaviour as \textbf{key-value pairs}.
  \item Values can be data (\textbf{properties}) or functions (\textbf{methods}).
  \item \textbf{JSON} (JavaScript Object Notation) is a text format for storing and exchanging object data.
  \item Technically, in JavaScript, almost everything is an object.
\end{itemize}

\begin{lstlisting}[language=Java, caption={Objects and JSON}]
// Object literal
const user = {
  name: "Alice",
  age: 25,
  greet() {
    return `Hi, I'm ${this.name}`;
  },
};

// Property access
user.name;          // dot notation
user["age"];        // bracket notation (useful for dynamic keys)

// Destructuring
const { name, age } = user;

// Spread operator
const updated = { ...user, age: 26 };

// JSON serialisation/deserialisation
const json = JSON.stringify(user);   // '{"name":"Alice","age":25}'
const parsed = JSON.parse(json);     // back to object
\end{lstlisting}

\subsection{Arrays}
\begin{lstlisting}[language=Java, caption={Array methods}]
const nums = [1, 2, 3, 4, 5];

nums.map(n => n * 2);          // [2, 4, 6, 8, 10] -- transform each element
nums.filter(n => n > 2);       // [3, 4, 5] -- keep matching elements
nums.reduce((acc, n) => acc + n, 0); // 15 -- fold to single value
nums.find(n => n > 3);         // 4 -- first match
nums.findIndex(n => n > 3);    // 3 -- index of first match
nums.some(n => n > 4);         // true -- at least one matches
nums.every(n => n > 0);        // true -- all match
nums.includes(3);              // true
nums.slice(1, 3);              // [2, 3] -- non-mutating
nums.splice(1, 2);             // removes & returns [2, 3], mutates original

// Spread and destructuring
const [first, second, ...rest] = nums;
const combined = [...nums, 6, 7];
\end{lstlisting}

\subsection{The \texttt{this} Keyword}
\begin{itemize}
  \item Refers to the object that is \textbf{calling the function}.
  \item In \textbf{objects}, \texttt{this} points to the object itself.
  \item In \textbf{classes}, \texttt{this} points to the created instance.
  \item In \textbf{arrow functions}, \texttt{this} is \textbf{inherited} from the outer (lexical) scope --- arrow functions have no own \texttt{this}.
  \item The value of \texttt{this} depends on \textbf{how} a function is called, not where it is defined.
\end{itemize}

\begin{lstlisting}[language=Java, caption={this in different contexts}]
const obj = {
  name: "Alice",
  greet() {
    console.log(this.name); // "Alice" -- this = obj
  },
  greetArrow: () => {
    console.log(this.name); // undefined -- arrow has no own this
  },
};

class Counter {
  count = 0;
  increment() {
    this.count++; // this = instance
  }
}
\end{lstlisting}

\subsection{Closures}
A \textbf{closure} is a function that \textbf{remembers} the variables from its outer scope even after that scope has finished executing.

\begin{lstlisting}[language=Java, caption={Closure example}]
function makeCounter() {
  let count = 0;        // outer variable
  return function() {   // inner function closes over `count`
    count++;
    return count;
  };
}

const counter = makeCounter();
counter(); // 1
counter(); // 2
counter(); // 3
// `count` is private and persists between calls
\end{lstlisting}

\subsection{The Event Loop and Asynchronous JavaScript}
JavaScript is \textbf{single-threaded} but uses an event loop to handle asynchronous operations without blocking.

\begin{itemize}
  \item \textbf{Call stack}: where currently executing code runs (LIFO).
  \item \textbf{Web APIs}: browser-provided APIs (timers, fetch, DOM events) that run outside the stack.
  \item \textbf{Callback queue} (task queue): completed async callbacks waiting to be run.
  \item \textbf{Microtask queue}: higher-priority queue for Promise callbacks (\texttt{.then}, \texttt{.catch}).
  \item The event loop pushes tasks from the queues to the stack only when the stack is \textbf{empty}.
  \item Microtasks always drain before the next macrotask runs.
\end{itemize}

\subsection{Callbacks}
\begin{itemize}
  \item A \textbf{callback} is a function passed as an argument to another function, to be called later.
  \item Simple and flexible, but can lead to deeply nested \textbf{``callback hell''}.
\end{itemize}

\begin{lstlisting}[language=Java, caption={Callback example and callback hell}]
// Simple callback
setTimeout(() => console.log("Done!"), 1000);

// Callback hell (pyramid of doom)
fetchUser(id, (user) => {
  fetchPosts(user.id, (posts) => {
    fetchComments(posts[0].id, (comments) => {
      console.log(comments); // deeply nested!
    });
  });
});
\end{lstlisting}

\subsection{Promises}
\begin{itemize}
  \item A \textbf{Promise} represents a value that will be available in the future.
  \item Three states: \textbf{pending}, \textbf{fulfilled}, \textbf{rejected}.
  \item Chain with \texttt{.then()} (success) and \texttt{.catch()} (error) and \texttt{.finally()}.
  \item \texttt{Promise.all()} --- waits for \textbf{all} promises; fails fast on any rejection.
  \item \texttt{Promise.allSettled()} --- waits for all, regardless of outcome.
  \item \texttt{Promise.race()} --- resolves/rejects as soon as the first settles.
\end{itemize}

\begin{lstlisting}[language=Java, caption={Promises}]
const fetchData = () =>
  new Promise((resolve, reject) => {
    setTimeout(() => resolve({ id: 1, name: "Alice" }), 500);
  });

fetchData()
  .then(user => { console.log(user.name); return user.id; })
  .then(id => console.log("ID:", id))
  .catch(err => console.error(err))
  .finally(() => console.log("Done"));

// Parallel execution
Promise.all([fetch("/users"), fetch("/posts")])
  .then(([users, posts]) => console.log(users, posts));
\end{lstlisting}

\subsection{Async / Await}
\begin{itemize}
  \item Built on top of Promises but with \textbf{synchronous-looking} syntax.
  \item \texttt{async} before a function makes it return a Promise automatically.
  \item \texttt{await} pauses execution inside the async function until the Promise resolves.
  \item Use \texttt{try/catch} for error handling.
  \item Easier to read than promise chains; preferred in modern code.
\end{itemize}

\begin{lstlisting}[language=Java, caption={async/await}]
async function loadUser(id) {
  try {
    const res = await fetch(`/api/users/${id}`);
    if (!res.ok) throw new Error("User not found");
    const user = await res.json();
    console.log(user.name);
    return user;
  } catch (err) {
    console.error("Failed:", err.message);
  }
}

// Parallel with async/await
async function loadAll() {
  const [users, posts] = await Promise.all([
    fetch("/users").then(r => r.json()),
    fetch("/posts").then(r => r.json()),
  ]);
  return { users, posts };
}
\end{lstlisting}

\subsection{ES Modules}
\begin{itemize}
  \item A \textbf{module} is a file with its own private scope.
  \item Use \texttt{export} to make values/functions visible to other files.
  \item Use \texttt{import} to use code from other files.
  \item \textbf{Named exports}: \texttt{export function foo() \{\}} --- imported as \texttt{\{ foo \}}.
  \item \textbf{Default export}: \texttt{export default ...} --- imported without braces.
  \item Modern browsers and Node.js support ES modules natively.
\end{itemize}

\begin{lstlisting}[language=Java, caption={ES module syntax}]
// utils.js
export const PI = 3.14159;
export function square(x) { return x * x; }
export default class MathHelper { /* ... */ }

// main.js
import MathHelper, { PI, square } from "./utils.js";
import * as Utils from "./utils.js"; // import all as namespace
\end{lstlisting}

\newpage
% ==================================================
\section{JavaScript Beyond the Browser}
% ==================================================

\subsection{Limitations of Plain JavaScript}
\begin{itemize}
  \item Originally designed for simple browser interactions, not full-scale applications.
  \item Single-threaded execution causes blocking during heavy I/O or computation.
  \item Dynamic typing leads to unpredictable runtime errors in large projects.
  \item Early JS lacked a proper module system, making code organisation difficult.
  \item Limited tooling and poor scalability for complex applications.
\end{itemize}

\subsection{Node.js}
\begin{itemize}
  \item Powered by Google's \textbf{V8 engine} --- the same engine that runs Chrome.
  \item \textbf{Event-driven, non-blocking I/O} model --- handles many requests concurrently.
  \item Same language (\textbf{JavaScript}) for both frontend and backend.
  \item Access to the \textbf{npm} ecosystem (millions of packages).
  \item Turned JavaScript into a complete server-side platform.
  \item Supports modern ES modules, \texttt{async/await}, and integrates with containers and cloud.
\end{itemize}

\subsection{NPM (Node Package Manager)}
\begin{itemize}
  \item Central registry to manage and distribute JavaScript packages.
  \item Standardises project setup with \texttt{package.json}.
  \item Key commands:
    \begin{itemize}
      \item \texttt{npm init} --- create a new project and generate \texttt{package.json}.
      \item \texttt{npm install <package>} --- add a dependency.
      \item \texttt{npm uninstall <package>} --- remove a package.
      \item \texttt{npm update} --- update all dependencies.
      \item \texttt{npm run <script>} --- execute a custom script defined in \texttt{package.json}.
    \end{itemize}
\end{itemize}

\begin{lstlisting}[language={}, caption={package.json scripts section}]
{
  "scripts": {
    "start":  "node server.js",
    "dev":    "nodemon server.js",
    "build":  "tsc",
    "test":   "jest"
  }
}
\end{lstlisting}

\subsection{Popular npm Packages}
\begin{itemize}
  \item \textbf{Express} --- minimalist web framework for APIs and servers.
  \item \textbf{React} --- library for building interactive UIs.
  \item \textbf{Lodash} --- utility functions for arrays, objects, strings.
  \item \textbf{Axios} --- promise-based HTTP client.
  \item \textbf{Socket.io} --- real-time bidirectional communication.
  \item \textbf{Jest} --- testing framework for JavaScript/Node.
  \item \textbf{Day.js} --- lightweight date/time manipulation.
  \item \textbf{Three.js} --- 3D graphics in the browser.
  \item \textbf{Electron} --- build cross-platform desktop apps with web technologies.
\end{itemize}

\subsection{Bundling for Browsers}
\begin{itemize}
  \item Combines multiple JS files into a single (or few) output files for browsers.
  \item Reduces HTTP requests and improves load performance.
  \item Handles module resolution, imports, and exports automatically.
  \item Often includes \textbf{minification} (removing whitespace) and \textbf{tree-shaking} (removing unused code).
  \item Popular tools: \textbf{Vite}, \textbf{Webpack}, \textbf{Rollup}, \textbf{Parcel}.
\end{itemize}

\subsection{TypeScript}
\subsubsection{Motivation}
\begin{itemize}
  \item JavaScript is \textbf{dynamically typed} --- variables can change type at runtime.
  \item Type coercion causes surprises: \texttt{"5" + 2} evaluates to \texttt{"52"}.
  \item Large codebases are harder to maintain due to hidden type bugs.
  \item No compile-time checks --- many errors appear only during execution.
\end{itemize}

\subsubsection{What is TypeScript?}
\begin{itemize}
  \item Created by \textbf{Microsoft} in 2012 (lead architect: \textbf{Anders Hejlsberg}).
  \item A \textbf{statically typed superset} of JavaScript --- all valid JS is valid TS.
  \item You write TypeScript $\to$ it compiles to JavaScript $\to$ runs anywhere JS runs.
  \item Designed for gradual, optional adoption.
  \item Widely used in enterprise and complex web applications.
\end{itemize}

\subsubsection{Variables and Types}
\begin{lstlisting}[language=Java, caption={TypeScript variable annotations}]
let age: number = 25;
const name: string = "Alice";
let isAdmin: boolean = false;

// Union type
let id: number | string = 42;
id = "abc-123"; // also valid

// TypeScript infers the type if not declared
let score = 100; // inferred as number
\end{lstlisting}

\subsubsection{Interfaces}
\begin{lstlisting}[language=Java, caption={TypeScript interface}]
interface User {
  id: number;
  name: string;
  email?: string;   // optional
  readonly role: string; // cannot be changed after creation
}

function greet(user: User): string {
  return `Hello, ${user.name}`;
}
\end{lstlisting}

\subsubsection{Classes}
\begin{lstlisting}[language=Java, caption={TypeScript class with access modifiers}]
class Animal {
  private name: string;
  protected age: number;

  constructor(name: string, age: number) {
    this.name = name;
    this.age = age;
  }

  public speak(): void {
    console.log(`${this.name} makes a sound.`);
  }
}

class Dog extends Animal {
  constructor(name: string, age: number) {
    super(name, age);
  }

  public speak(): void {
    console.log("Woof!");
  }
}
\end{lstlisting}

\subsubsection{Generics}
\begin{lstlisting}[language=Java, caption={TypeScript generics}]
function identity<T>(arg: T): T {
  return arg;
}

identity<string>("hello"); // returns "hello"
identity<number>(42);      // returns 42

// Generic interface
interface ApiResponse<T> {
  data: T;
  status: number;
  message: string;
}
\end{lstlisting}

\subsubsection{Modules}
\begin{lstlisting}[language=Java, caption={ES module exports and imports in TypeScript}]
// math.ts
export function add(a: number, b: number): number {
  return a + b;
}

export default function multiply(a: number, b: number): number {
  return a * b;
}

// main.ts
import multiply, { add } from "./math";
console.log(add(2, 3));       // 5
console.log(multiply(2, 3));  // 6
\end{lstlisting}

\newpage
% ==================================================
\section{The DOM and Browser APIs}
% ==================================================

\subsection{From Markup to DOM}
\begin{itemize}
  \item The browser \textbf{parses HTML} into a tree of nodes --- the \textbf{DOM (Document Object Model)}.
  \item Each HTML tag becomes an \textbf{element node}.
  \item Text inside tags becomes \textbf{text nodes}.
  \item Attributes become \textbf{node properties}.
  \item The DOM is what JavaScript can access and modify.
\end{itemize}

\subsection{DOM Structure}
\begin{itemize}
  \item \textbf{Root node}: \texttt{document}.
  \item \textbf{window} --- top-level browser object (global scope).
  \item \textbf{document} --- represents the web page inside the window.
  \item \textbf{Elements} --- tags (\texttt{<div>}, \texttt{<p>}) as objects in the tree.
  \item \textbf{Attributes \& Text} --- properties of element nodes.
  \item \textbf{Events} --- methods and listeners that connect JS to user actions.
\end{itemize}

\subsection{Querying DOM Elements}
\begin{lstlisting}[language=Java, caption={DOM querying methods}]
// By ID (returns a single Element or null)
const btn = document.getElementById("submit-btn");

// By class name (returns HTMLCollection)
const cards = document.getElementsByClassName("card");

// By tag name (returns HTMLCollection)
const paragraphs = document.getElementsByTagName("p");

// CSS selector -- first match (most flexible)
const nav = document.querySelector("header > nav");

// CSS selector -- all matches (returns NodeList)
const links = document.querySelectorAll("a[href^='https']");
\end{lstlisting}

\subsection{Events}
\begin{itemize}
  \item Events let webpages react to user actions.
  \item Common events: \texttt{click}, \texttt{input}, \texttt{change}, \texttt{submit}, \texttt{keydown}, \texttt{keyup}, \texttt{load}, \texttt{scroll}, \texttt{resize}.
  \item \texttt{addEventListener} is the \textbf{modern, preferred} way to attach events.
  \item Multiple handlers can be added to the same event on the same element.
\end{itemize}

\begin{lstlisting}[language=Java, caption={Event handling example}]
const button = document.querySelector("#myBtn");

button.addEventListener("click", (event) => {
  console.log("Clicked!", event.target);
  event.preventDefault(); // prevent default browser action
});

// Remove listener
button.removeEventListener("click", handlerFunction);
\end{lstlisting}

\subsection{Manipulating the DOM}

\subsubsection{Modifying Content and Styles}
\begin{lstlisting}[language=Java, caption={Changing content and styles}]
const heading = document.querySelector("h1");

heading.textContent = "New Title";         // safe (plain text)
heading.innerHTML = "<em>New</em> Title";  // parses HTML

heading.style.color = "crimson";
heading.style.fontSize = "2rem";

heading.setAttribute("data-active", "true");
heading.getAttribute("data-active"); // "true"
heading.removeAttribute("data-active");
\end{lstlisting}

\subsubsection{Adding and Removing Elements}
\begin{lstlisting}[language=Java, caption={Creating and removing DOM nodes}]
// Create and insert
const div = document.createElement("div");
div.classList.add("card");
div.textContent = "Hello!";
document.body.appendChild(div);

// Insert before another element
const list = document.querySelector("ul");
const newItem = document.createElement("li");
newItem.textContent = "New item";
list.insertBefore(newItem, list.firstChild);

// Remove
div.remove();
// or: list.removeChild(newItem);
\end{lstlisting}

\subsubsection{Working with Classes}
\begin{lstlisting}[language=Java, caption={classList API}]
element.classList.add("active");
element.classList.remove("hidden");
element.classList.toggle("dark-mode");     // adds if absent, removes if present
element.classList.contains("active");      // returns boolean
element.classList.replace("old", "new");
\end{lstlisting}

\subsection{Forms and DOM}
\begin{lstlisting}[language=Java, caption={Reading form input values}]
const form = document.querySelector("#loginForm");
const emailInput = document.querySelector("#email");

form.addEventListener("submit", (e) => {
  e.preventDefault(); // stop page reload
  const email = emailInput.value;
  const isChecked = document.querySelector("#remember").checked;
  console.log(email, isChecked);
});
\end{lstlisting}

\subsection{The \texttt{window} Object}
\begin{itemize}
  \item Represents the browser window; all global variables and functions belong to it.
  \item \texttt{alert()}, \texttt{confirm()}, \texttt{prompt()} --- dialogs.
  \item \texttt{location} --- current URL; can be used to redirect: \texttt{window.location.href = "/new"}.
  \item \texttt{history} --- browser history (\texttt{back()}, \texttt{forward()}).
  \item \texttt{navigator} --- browser/device information.
  \item \texttt{setTimeout(fn, ms)} / \texttt{setInterval(fn, ms)} --- timers.
  \item \texttt{innerWidth} / \texttt{innerHeight} --- viewport dimensions.
  \item \texttt{scrollX} / \texttt{scrollY} --- scroll position.
\end{itemize}

\subsection{Web Storage}
\subsubsection{\texttt{localStorage}}
\begin{itemize}
  \item Stores data \textbf{persistently} in the browser (survives closing the tab/browser).
  \item Data stored as \textbf{key-value pairs of strings}.
  \item Same-origin policy: data is isolated per domain.
  \item Approximately 5--10 MB per site.
  \item \textbf{Synchronous} and not safe for sensitive information.
\end{itemize}
\begin{lstlisting}[language=Java, caption={localStorage usage}]
// Store
localStorage.setItem("theme", "dark");
localStorage.setItem("user", JSON.stringify({ name: "Alice" }));

// Retrieve
const theme = localStorage.getItem("theme"); // "dark"
const user = JSON.parse(localStorage.getItem("user"));

// Remove
localStorage.removeItem("theme");
localStorage.clear(); // removes all keys
\end{lstlisting}

\subsubsection{\texttt{sessionStorage}}
\begin{itemize}
  \item Data is cleared when the \textbf{tab or window is closed}.
  \item Same API as \texttt{localStorage}: \texttt{setItem}, \texttt{getItem}, \texttt{removeItem}, \texttt{clear}.
  \item Useful for temporary session data (e.g.\ multi-step form progress).
\end{itemize}

\newpage
% ==================================================
\section{React --- UIs That React}
% ==================================================

\subsection{Motivation: Problems with Classic Frontend Development}
\begin{itemize}
  \item Tangled DOM manipulation with vanilla JS / jQuery.
  \item Inconsistent and hard-to-manage \textbf{state}.
  \item Poor code reuse and maintainability.
  \item Performance bottlenecks from manual DOM updates.
  \item Difficult collaboration and scaling.
\end{itemize}

\subsection{The Rise and Limits of jQuery}
\begin{itemize}
  \item jQuery made DOM manipulation simpler and faster, and abstracted cross-browser inconsistencies.
  \item Popularised easy AJAX requests and brought a cleaner API than raw JavaScript.
  \item However, jQuery still left \textbf{state management} and \textbf{application structure} problems unsolved.
  \item As applications grew more complex, jQuery codebases became difficult to maintain and reason about.
  \item This led to the rise of modern frameworks and libraries like Angular, Vue, and React.
\end{itemize}

\subsection{Library vs.\ Framework}
\begin{itemize}
  \item \textbf{Library}: you call it --- you are in control (e.g.\ React, Axios).
  \item \textbf{Framework}: it calls you --- it is in control (e.g.\ Angular, Next.js).
  \item Libraries provide tools; frameworks provide structure and enforce conventions.
  \item \textbf{React is a library}, not a framework --- it handles the view layer and leaves routing, state, and data-fetching to other tools.
\end{itemize}

\subsection{What is React?}
\begin{itemize}
  \item A \textbf{JavaScript library} for building user interfaces (not a full framework).
  \item Created by \textbf{Facebook} in 2013, now maintained by Meta and the open-source community.
  \item \textbf{Component-based}: build UIs from small, reusable pieces.
  \item \textbf{Declarative}: describe \emph{what} the UI should look like, not \emph{how} to update it.
  \item Efficient rendering via the \textbf{Virtual DOM}.
\end{itemize}

\subsection{Setting Up a React Project with Vite}
\begin{lstlisting}[language={}, caption={Creating a React + Vite project}]
npm create vite@latest my-app
# Select "React" or "React + TypeScript" when prompted

cd my-app
npm install
npm run dev
# Opens at http://localhost:5173
\end{lstlisting}

\begin{itemize}
  \item \textbf{Vite} uses native ES modules in development, providing near-instant HMR (Hot Module Replacement).
  \item Uses \textbf{Rollup} for optimised production builds.
  \item Created by Evan You (Vue.js creator).
\end{itemize}

\subsection{JSX}
\begin{itemize}
  \item A \textbf{syntax extension} for JavaScript that lets you write HTML-like code inside JS.
  \item \textbf{Not valid JS} --- compiled to plain JavaScript by \textbf{Babel}.
  \item Rules:
    \begin{itemize}
      \item Must return a \textbf{single root element} (use \texttt{<>...</>} fragments if needed).
      \item Attribute names use \textbf{camelCase}: \texttt{className}, \texttt{onClick}, \texttt{htmlFor}.
      \item Self-closing tags must close: \texttt{<img />}, \texttt{<input />}.
      \item JavaScript \textbf{expressions} go inside \texttt{\{\}}.
      \item Statements (\texttt{if}, \texttt{for}) are \textbf{not} allowed directly inside JSX.
    \end{itemize}
\end{itemize}

\begin{lstlisting}[language=Java, caption={JSX examples}]
const name = "World";
const isLoggedIn = true;

return (
  <div className="container">
    <h1>Hello, {name}!</h1>
    <p>{isLoggedIn ? "Welcome back." : "Please log in."}</p>
    <ul>
      {["Apple", "Banana", "Cherry"].map((fruit, i) => (
        <li key={i}>{fruit}</li>
      ))}
    </ul>
  </div>
);
\end{lstlisting}

\subsection{React Components}
\begin{itemize}
  \item Independent, reusable building blocks of the UI.
  \item Accept \textbf{props} as input and return JSX.
  \item Two types: \textbf{Function-based} (modern, preferred) and Class-based (legacy).
  \item Component names must start with an \textbf{uppercase letter}.
\end{itemize}

\begin{lstlisting}[language=Java, caption={A functional React component}]
interface ButtonProps {
  label: string;
  onClick: () => void;
  disabled?: boolean;
}

function Button({ label, onClick, disabled = false }: ButtonProps) {
  return (
    <button onClick={onClick} disabled={disabled} className="btn">
      {label}
    </button>
  );
}

export default Button;
\end{lstlisting}

\subsection{Props}
\begin{itemize}
  \item \textbf{Props} are inputs passed from parent to child components.
  \item \textbf{Read-only}: child components must not modify them.
  \item Passed like HTML attributes.
  \item The special \texttt{children} prop passes nested JSX into a component.
\end{itemize}

\begin{lstlisting}[language=Java, caption={Using props and children}]
function Card({ title, children }: { title: string; children: React.ReactNode }) {
  return (
    <div className="card">
      <h2>{title}</h2>
      <div className="card-body">{children}</div>
    </div>
  );
}

// Usage:
<Card title="Hello">
  <p>This is the card body content.</p>
</Card>
\end{lstlisting}

\subsection{State and \texttt{useState}}
\begin{itemize}
  \item \textbf{State} is component memory that stores dynamic data.
  \item Changing state \textbf{triggers a re-render}.
  \item Managed with the \texttt{useState} hook: \texttt{const [value, setValue] = useState(initialValue)}.
  \item Always update via the setter --- \textbf{never modify state directly}.
  \item Use a \textbf{callback form} when new state depends on old state: \texttt{setValue(v => v + 1)}.
\end{itemize}

\begin{lstlisting}[language=Java, caption={useState example --- counter}]
import { useState } from "react";

function Counter() {
  const [count, setCount] = useState(0);

  return (
    <div>
      <p>Count: {count}</p>
      <button onClick={() => setCount(c => c + 1)}>+</button>
      <button onClick={() => setCount(c => c - 1)}>-</button>
      <button onClick={() => setCount(0)}>Reset</button>
    </div>
  );
}
\end{lstlisting}

\subsection{Event Handling in React}
\begin{itemize}
  \item Use camelCase event names: \texttt{onClick}, \texttt{onChange}, \texttt{onSubmit}.
  \item Pass a \textbf{function reference}, not a function call: \texttt{onClick=\{handleClick\}} not \texttt{onClick=\{handleClick()\}}.
  \item Access the event object: \texttt{onClick=\{(e) => console.log(e.target)\}}.
  \item Prevent default behaviour: \texttt{e.preventDefault()}.
\end{itemize}

\subsection{Lists in JSX}
\begin{lstlisting}[language=Java, caption={Rendering a list with keys}]
interface Post { id: number; title: string; }

function PostList({ posts }: { posts: Post[] }) {
  return (
    <ul>
      {posts.map((post) => (
        <li key={post.id}>{post.title}</li>
      ))}
    </ul>
  );
}
\end{lstlisting}
Always provide a stable, unique \texttt{key} for list items. Avoid using array index as key when the list can reorder.

\subsection{The Virtual DOM and Reconciliation}
\begin{itemize}
  \item React maintains a \textbf{Virtual DOM} --- a lightweight in-memory copy of the real DOM.
  \item When state or props change, React builds a \textbf{new VDOM tree}.
  \item The \textbf{Fiber reconciler} diffs the new tree against the old one (\textbf{O(n)} algorithm with heuristics).
  \item Only the \textbf{minimal set of real DOM updates} is applied.
  \item Heuristics: (1) Different element types always produce new subtrees; (2) \texttt{key} props help track moved elements.
\end{itemize}

\subsection{Component Composition}
\begin{itemize}
  \item Building UIs by \textbf{combining} smaller components together.
  \item Promotes reusability and clean separation of concerns.
  \item Components can nest inside others: \texttt{<Layout><Header /></Layout>}
  \item Use \texttt{props} and \texttt{children} to pass content into components.
  \item Encourages modular, maintainable React architecture.
  \item Prefer \textbf{composition} over inheritance for sharing logic between components.
\end{itemize}

\subsection{React Developer Tools}
\begin{itemize}
  \item \textbf{React DevTools} is a browser extension for inspecting React applications.
  \item Allows inspection of the \textbf{component tree} and current props and state.
  \item Highlights re-renders to help identify performance issues.
  \item Enables \textbf{editing of props and state} during development without code changes.
  \item Useful for debugging hooks, context, and component hierarchy.
  \item Also use \texttt{console.log()} in key spots to trace state and prop values.
\end{itemize}

\subsection{Styling in React}
\begin{itemize}
  \item \textbf{Plain CSS}: import \texttt{.css} files; apply with \texttt{className}.
  \item \textbf{CSS Modules}: scoped styles per component (filename: \texttt{Button.module.css}).
  \item \textbf{Inline styles}: \texttt{style=\{\{ backgroundColor: "red", fontSize: "1rem" \}\}} (camelCase).
  \item \textbf{TailwindCSS}: utility-first classes applied directly in JSX.
  \item \textbf{shadcn/ui}: accessible, copyable React components built on Radix UI + Tailwind.
\end{itemize}

\newpage
% ==================================================
\section{React --- Act II: Hooks, Lifecycle, and Patterns}
% ==================================================

\subsection{Component Lifecycle}
\begin{itemize}
  \item \textbf{Mount}: component is created and inserted into the DOM.
  \item \textbf{Update}: state or props change, triggering a re-render.
  \item \textbf{Unmount}: component is removed from the DOM.
  \item Rendering must remain \textbf{pure} --- same inputs produce the same output; no side effects during render.
  \item Side effects are managed with \texttt{useEffect}.
\end{itemize}

\subsection{The \texttt{useEffect} Hook}
\begin{itemize}
  \item Runs \textbf{after} the component renders.
  \item Used to synchronise the component with external systems: APIs, timers, subscriptions, DOM mutations.
  \item Accepts a \textbf{dependency array} that controls when the effect re-runs.
  \item Can return a \textbf{cleanup function} to undo the effect on unmount or before re-running.
\end{itemize}

\begin{lstlisting}[language=Java, caption={useEffect variations}]
// Run once on mount
useEffect(() => {
  console.log("Component mounted");
  return () => console.log("Component unmounted"); // cleanup
}, []);

// Run when `userId` changes
useEffect(() => {
  fetchUser(userId).then(setUser);
}, [userId]);

// Run after every render (omit dependency array -- rarely needed)
useEffect(() => {
  document.title = `Count: ${count}`;
});
\end{lstlisting}

\subsubsection{Common \texttt{useEffect} Mistakes}
\begin{itemize}
  \item \textbf{Missing dependencies}: causes stale closures and hard-to-debug bugs.
  \item \textbf{Calling \texttt{setState} unconditionally inside effects}: can cause infinite re-render loops.
  \item \textbf{Deriving state from props inside effects}: compute it during render instead.
  \item Effects may run \textbf{twice in Strict Mode} (React mounts, unmounts, then mounts again to detect side effects).
\end{itemize}

\subsection{Fetching Data}

\subsubsection{Using \texttt{fetch}}
\begin{lstlisting}[language=Java, caption={Data fetching with fetch + useEffect}]
import { useState, useEffect } from "react";

function UserList() {
  const [users, setUsers] = useState([]);
  const [loading, setLoading] = useState(true);
  const [error, setError] = useState(null);

  useEffect(() => {
    let cancelled = false;
    setLoading(true);
    fetch("https://api.example.com/users")
      .then((res) => {
        if (!res.ok) throw new Error("Network error");
        return res.json();
      })
      .then((data) => { if (!cancelled) setUsers(data); })
      .catch((err) => { if (!cancelled) setError(err.message); })
      .finally(() => { if (!cancelled) setLoading(false); });

    return () => { cancelled = true; }; // prevent state update after unmount
  }, []);

  if (loading) return <p>Loading...</p>;
  if (error)   return <p>Error: {error}</p>;
  return <ul>{users.map(u => <li key={u.id}>{u.name}</li>)}</ul>;
}
\end{lstlisting}

\subsubsection{Using Axios}
\begin{lstlisting}[language=Java, caption={Data fetching with Axios}]
import axios from "axios";
import { useState, useEffect } from "react";

function UserList() {
  const [users, setUsers] = useState([]);

  useEffect(() => {
    const source = axios.CancelToken.source();
    axios.get("/api/users", { cancelToken: source.token })
      .then(({ data }) => setUsers(data))
      .catch((err) => {
        if (!axios.isCancel(err)) console.error(err);
      });
    return () => source.cancel("Component unmounted");
  }, []);

  return <ul>{users.map(u => <li key={u.id}>{u.name}</li>)}</ul>;
}
\end{lstlisting}

\subsection{State Management Patterns}

\subsubsection{Pattern I: Lifting State Up}
\begin{itemize}
  \item Move shared state to the \textbf{closest common ancestor}.
  \item Children receive data as props; updates flow upward via callbacks.
  \item Prevents state divergence between siblings.
\end{itemize}

\subsubsection{Pattern II: Callbacks as Props}
\begin{itemize}
  \item Parent passes a function to a child for state updates.
  \item The child calls the function to communicate intent, without owning state.
  \item Enables \textbf{unidirectional data flow}.
  \item May lead to \textbf{prop drilling} in deep component trees.
\end{itemize}

\begin{lstlisting}[language=Java, caption={Callback as prop}]
function Parent() {
  const [value, setValue] = useState("");
  return <Input value={value} onChange={setValue} />;
}

function Input({ value, onChange }: { value: string; onChange: (v: string) => void }) {
  return (
    <input value={value} onChange={(e) => onChange(e.target.value)} />
  );
}
\end{lstlisting}

\subsubsection{Pattern III: Context API}
\begin{itemize}
  \item Provides a way to share state \textbf{without prop drilling}.
  \item Data is supplied via a \texttt{Provider} and consumed by any descendant via \texttt{useContext}.
  \item Best for \textbf{global concerns}: theme, authentication, locale, feature flags.
  \item \textbf{Overusing Context} can cause unnecessary re-renders --- scope it tightly.
\end{itemize}

\begin{lstlisting}[language=Java, caption={Context API example}]
import { createContext, useContext, useState } from "react";

const ThemeContext = createContext<"dark" | "light">("dark");

function App() {
  const [theme, setTheme] = useState<"dark" | "light">("dark");
  return (
    <ThemeContext.Provider value={theme}>
      <button onClick={() => setTheme(t => t === "dark" ? "light" : "dark")}>
        Toggle
      </button>
      <Page />
    </ThemeContext.Provider>
  );
}

function Page() {
  const theme = useContext(ThemeContext);
  return <div className={theme}>Page content</div>;
}
\end{lstlisting}

\subsection{Redux}
\begin{itemize}
  \item Predictable state management library for React.
  \item Application state lives in a \textbf{single global store}.
  \item State updates occur through \textbf{dispatched actions} and \textbf{pure reducer} functions.
  \item Enables centralised debugging and \textbf{time-travel inspection}.
  \item Best suited for large applications with complex shared state.
\end{itemize}

\subsection{React Router}
\begin{itemize}
  \item Enables \textbf{client-side routing} in SPAs --- maps URLs to components without full page reloads.
  \item Provided by the \texttt{react-router-dom} package.
  \item Core components: \texttt{BrowserRouter}, \texttt{Routes}, \texttt{Route}.
  \item Hooks: \texttt{useParams}, \texttt{useNavigate}, \texttt{useLocation}.
  \item Supports nested routes, dynamic segments, and layout composition.
\end{itemize}

\begin{lstlisting}[language=Java, caption={React Router v6 example}]
import { BrowserRouter, Routes, Route, Link } from "react-router-dom";

function App() {
  return (
    <BrowserRouter>
      <nav>
        <Link to="/">Home</Link>
        <Link to="/about">About</Link>
        <Link to="/users/42">User 42</Link>
      </nav>
      <Routes>
        <Route path="/"          element={<Home />} />
        <Route path="/about"     element={<About />} />
        <Route path="/users/:id" element={<UserDetail />} />
        <Route path="*"          element={<NotFound />} />
      </Routes>
    </BrowserRouter>
  );
}
\end{lstlisting}

\subsection{Forms in React}
\begin{itemize}
  \item \textbf{Controlled components}: form input values are driven by React state; use \texttt{value} + \texttt{onChange}.
  \item \textbf{Uncontrolled components}: rely on \texttt{useRef} to access DOM values directly --- simpler but less React-idiomatic.
  \item Form submission handled via event handlers rather than default browser behaviour.
  \item Always call \texttt{e.preventDefault()} to stop the page from reloading on submit.
  \item Validation can be synchronous (in state) or asynchronous (with effects or libraries like \textbf{React Hook Form}).
\end{itemize}

\begin{lstlisting}[language=Java, caption={Controlled form in React}]
import { useState } from "react";

function LoginForm() {
  const [email, setEmail] = useState("");
  const [password, setPassword] = useState("");
  const [error, setError] = useState("");

  const handleSubmit = (e: React.FormEvent) => {
    e.preventDefault();
    if (!email.includes("@")) {
      setError("Invalid email address");
      return;
    }
    console.log("Submitting:", { email, password });
  };

  return (
    <form onSubmit={handleSubmit}>
      <input
        type="email"
        value={email}
        onChange={(e) => setEmail(e.target.value)}
        placeholder="Email"
      />
      <input
        type="password"
        value={password}
        onChange={(e) => setPassword(e.target.value)}
        placeholder="Password"
      />
      {error && <p style={{ color: "red" }}>{error}</p>}
      <button type="submit">Log In</button>
    </form>
  );
}
\end{lstlisting}

\subsection{Error Handling in React}
\begin{itemize}
  \item Errors can occur during \textbf{rendering}, \textbf{data fetching}, or \textbf{user interactions}.
  \item \textbf{Error Boundaries}: class components that catch render-time errors in their subtree and display a fallback UI instead of crashing.
  \item Async errors (fetch failures, timeouts) must be caught explicitly in \texttt{useEffect} or event handlers.
  \item Runtime errors should always be anticipated and handled gracefully.
  \item Proper error handling improves resilience and user experience.
\end{itemize}

\begin{lstlisting}[language=Java, caption={Simple Error Boundary class component}]
import { Component, ReactNode } from "react";

interface State { hasError: boolean; }

class ErrorBoundary extends Component<{ children: ReactNode }, State> {
  state: State = { hasError: false };

  static getDerivedStateFromError() {
    return { hasError: true };
  }

  componentDidCatch(error: Error, info: any) {
    console.error("Caught:", error, info);
  }

  render() {
    if (this.state.hasError)
      return <h2>Something went wrong. Please refresh.</h2>;
    return this.props.children;
  }
}

// Usage
<ErrorBoundary>
  <MyComponent />
</ErrorBoundary>
\end{lstlisting}

\subsection{shadcn/ui}
\begin{itemize}
  \item A collection of \textbf{reusable, accessible} React components.
  \item Components are \textbf{copied into your project} (not installed as a library dependency) --- you own the code.
  \item Built on top of \textbf{Radix UI} (headless, accessible primitives) and \textbf{TailwindCSS}.
  \item Fully customisable since the source lives in your codebase.
  \item Encourages design consistency without locking into a fixed component library.
  \item Install individual components: \texttt{npx shadcn-ui@latest add button}.
\end{itemize}

\subsection{TanStack Query}
\begin{itemize}
  \item Manages \textbf{server/remote state} that can change outside the UI.
  \item Standardises loading, error, caching, and refetching behaviour.
  \item Reduces boilerplate compared to manual \texttt{useEffect} + \texttt{useState}.
  \item Prevents duplicate requests via \textbf{request deduplication}.
  \item Improves perceived performance via caching and background updates.
\end{itemize}

\begin{lstlisting}[language=Java, caption={TanStack Query --- useQuery example}]
import { useQuery, useMutation, useQueryClient } from "@tanstack/react-query";
import axios from "axios";

function Users() {
  const { data, isLoading, error } = useQuery({
    queryKey: ["users"],
    queryFn: () => axios.get("/api/users").then(r => r.data),
    staleTime: 60_000, // 1 min before background refetch
  });

  if (isLoading) return <p>Loading...</p>;
  if (error)     return <p>Error!</p>;

  return <ul>{data.map((u: any) => <li key={u.id}>{u.name}</li>)}</ul>;
}

function CreateUser() {
  const queryClient = useQueryClient();
  const mutation = useMutation({
    mutationFn: (newUser: any) => axios.post("/api/users", newUser),
    onSuccess: () => queryClient.invalidateQueries({ queryKey: ["users"] }),
  });

  return (
    <button onClick={() => mutation.mutate({ name: "New User" })}>
      Add User
    </button>
  );
}
\end{lstlisting}

\subsection{TanStack Mutations}
\begin{itemize}
  \item \texttt{useMutation} handles \textbf{create / update / delete} operations against the server.
  \item \texttt{onSuccess} handlers typically \textbf{invalidate related queries} to refresh cached data.
  \item \textbf{Optimistic updates}: update the UI before the server responds for a snappier feel.
  \item \textbf{Rollbacks}: revert the optimistic update if the server request fails.
  \item Mutation lifecycle callbacks: \texttt{onMutate}, \texttt{onSuccess}, \texttt{onError}, \texttt{onSettled}.
\end{itemize}

\subsection{TanStack Best Practices}
\begin{itemize}
  \item Treat TanStack Query as the default for \textbf{API/server state} --- use React state or Context for \textbf{UI state}.
  \item Normalise query keys and keep them consistent across the app.
  \item Prefer \textbf{invalidation} over manual refetch wiring for correctness.
  \item Use \texttt{select} / transform patterns to keep components minimal.
  \item Combine with lightweight client state (Context or Zustand) when needed.
\end{itemize}

\subsection{Next.js and SSR}
\begin{itemize}
  \item \textbf{SSR (Server-Side Rendering)}: HTML is rendered on the \textbf{server} and sent to the client.
  \item Improves \textbf{initial load performance} and \textbf{SEO} (content visible to crawlers immediately).
  \item Requires \textbf{hydration} to attach React interactivity on the client.
  \item \textbf{Next.js}: production-grade React framework with built-in SSR, SSG, file-based routing, and performance optimisations.
  \item Supports multiple rendering strategies: SSR, SSG (Static Site Generation), CSR (Client-Side Rendering).
\end{itemize}

\subsection{Recommended Project Structure}
\begin{lstlisting}[language={}, caption={Typical React project layout}]
src/
  assets/          # images, icons, fonts
  components/      # reusable UI components (Button, Card, Modal)
  pages/           # route-level components (HomePage, AboutPage)
  hooks/           # custom React hooks (useAuth, useFetch)
  services/        # API calls and external integrations
  context/         # React Context providers
  store/           # Redux / Zustand state
  utils/           # helper functions
  types/           # TypeScript interfaces and types
  App.tsx
  main.tsx
\end{lstlisting}

\newpage
% ==================================================
\section{Quick Reference}
% ==================================================

\subsection{JavaScript Cheat Sheet}
\begin{itemize}
  \item \textbf{Declarations}: \texttt{const} (preferred) $>$ \texttt{let} $>$ \texttt{var} (avoid).
  \item \textbf{Strict equality}: always use \texttt{===} / \texttt{!==}; avoid \texttt{==} to prevent coercion surprises.
  \item \textbf{Falsy values}: \texttt{false}, \texttt{0}, \texttt{""}, \texttt{null}, \texttt{undefined}, \texttt{NaN}.
  \item \textbf{Array methods}: \texttt{.map()}, \texttt{.filter()}, \texttt{.reduce()}, \texttt{.find()}, \texttt{.some()}, \texttt{.every()}.
  \item \textbf{Object spread}: \texttt{\{ ...obj, key: value \}} creates a shallow copy with overrides.
  \item \textbf{Optional chaining}: \texttt{user?.address?.city} --- short-circuits to \texttt{undefined} instead of throwing.
  \item \textbf{Nullish coalescing}: \texttt{value ?? "default"} --- uses right side only when left is \texttt{null}/\texttt{undefined}.
  \item \textbf{Async}: Callbacks $\to$ Promises $\to$ \texttt{async/await} (cleanest).
  \item \textbf{Modules}: \texttt{export} / \texttt{import}; default vs.\ named exports.
\end{itemize}

\subsection{HTML Cheat Sheet}
\begin{itemize}
  \item \texttt{<!DOCTYPE html>} --- HTML5 declaration.
  \item \texttt{<html>}, \texttt{<head>}, \texttt{<body>} --- document skeleton.
  \item \texttt{<h1>}--\texttt{<h6>}, \texttt{<p>}, \texttt{<strong>}, \texttt{<em>} --- text.
  \item \texttt{<a href>}, \texttt{<img src alt>}, \texttt{<video>}, \texttt{<audio>} --- links and media.
  \item \texttt{<ul>}, \texttt{<ol>}, \texttt{<li>}, \texttt{<dl>}, \texttt{<dt>}, \texttt{<dd>} --- lists.
  \item \texttt{<table>}, \texttt{<tr>}, \texttt{<th>}, \texttt{<td>} --- tables.
  \item \texttt{<form>}, \texttt{<input>}, \texttt{<select>}, \texttt{<textarea>}, \texttt{<button>} --- forms.
  \item \texttt{<header>}, \texttt{<nav>}, \texttt{<main>}, \texttt{<section>}, \texttt{<footer>} --- semantics.
\end{itemize}

\subsection{CSS Cheat Sheet}
\begin{itemize}
  \item Selectors: \texttt{element}, \texttt{.class}, \texttt{\#id}, \texttt{A > B}, \texttt{A + B}, \texttt{:hover}, \texttt{::before}.
  \item Box model: \texttt{margin} $\to$ \texttt{border} $\to$ \texttt{padding} $\to$ \texttt{content}.
  \item Layout: \texttt{display: flex | grid | block | inline}.
  \item Position: \texttt{static | relative | absolute | fixed | sticky}.
  \item Responsive: \texttt{@media (min-width: 768px) \{ ... \}}.
  \item Variables: \texttt{--name: value;} and \texttt{var(--name)}.
\end{itemize}

\subsection{React Hooks Summary}
\begin{itemize}
  \item \texttt{useState(init)} --- local component state; triggers re-render on change.
  \item \texttt{useEffect(fn, deps)} --- side effects (fetch, subscriptions, timers); runs after render.
  \item \texttt{useContext(Context)} --- consume a context value provided by a \texttt{Provider} ancestor.
  \item \texttt{useRef(init)} --- mutable ref that persists across renders without causing re-renders; also used for DOM access.
  \item \texttt{useMemo(fn, deps)} --- memoised computed value; re-computes only when dependencies change.
  \item \texttt{useCallback(fn, deps)} --- memoised function reference; prevents unnecessary child re-renders.
  \item \texttt{useReducer(reducer, init)} --- complex state logic with actions (like a mini Redux).
  \item \texttt{useLayoutEffect(fn, deps)} --- like \texttt{useEffect} but fires \textbf{synchronously} before the browser paints.
  \item \textbf{Rules of Hooks}: only call at the \textbf{top level} (not inside loops/conditions), only in \textbf{React function} components or custom hooks.
\end{itemize}

\end{document}